\documentclass[11pt,a4paper]{article}

% ===================== FONTS & LANGUAGE (XeLaTeX) =====================
\usepackage{fontspec}
\usepackage{polyglossia}
\setmainlanguage{vietnamese}
\setmainfont{Times New Roman}
\newfontfamily\headingfont{Times New Roman}[BoldFont={Times New Roman Bold}]
\newfontfamily\monofont{Courier New}

% ===================== LAYOUT =====================
\usepackage[a4paper,margin=2cm]{geometry}
\usepackage{pdflscape}
\usepackage{longtable}
\usepackage{array}
\usepackage{booktabs}
\usepackage{caption}
\usepackage{enumitem}
\usepackage{titlesec}
\usepackage[table]{xcolor}
\usepackage{hyperref}

\definecolor{primaryblue}{RGB}{31,73,125}
\definecolor{lightblue}{RGB}{221,235,247}
\definecolor{actorgray}{RGB}{90,90,90}
\definecolor{ucfill}{RGB}{255,242,204}
\definecolor{extendfill}{RGB}{252,213,180}
\definecolor{screenfill}{RGB}{198,224,180}
\definecolor{hubfill}{RGB}{180,199,231}

\hypersetup{
	colorlinks=true,
	linkcolor=primaryblue,
	urlcolor=primaryblue,
	citecolor=primaryblue,
	pdftitle={SRS - He thong ca nhan hoa lo trinh hoc va luyen thi danh gia nang luc tich hop AI},
	pdfauthor={SRS Document}
}

\titleformat{\section}{\headingfont\Large\bfseries\color{primaryblue}}{\thesection}{0.6em}{}
\titleformat{\subsection}{\headingfont\large\bfseries\color{primaryblue}}{\thesubsection}{0.6em}{}
\titleformat{\subsubsection}{\headingfont\normalsize\bfseries}{\thesubsubsection}{0.6em}{}

\setlength{\parindent}{0pt}
\setlength{\parskip}{6pt}

% ===================== TIKZ =====================
\usepackage{tikz}
\usetikzlibrary{shapes.geometric,shapes.symbols,arrows.meta,positioning,calc,backgrounds,fit,decorations.pathmorphing}

\tikzset{
	screennode/.style={
		rectangle, rounded corners=2pt, draw=primaryblue, thick,
		fill=screenfill, align=center, font=\footnotesize\sffamily,
		text width=2.5cm, minimum height=1.05cm, inner sep=2pt
	},
	hubnode/.style={
		rectangle, rounded corners=4pt, draw=primaryblue, very thick,
		fill=hubfill, align=center, font=\small\sffamily\bfseries,
		text width=2.6cm, minimum height=1.2cm, inner sep=3pt
	},
	loginnode/.style={
		rectangle, rounded corners=8pt, draw=primaryblue, very thick,
		fill=lightblue, align=center, font=\small\sffamily\bfseries,
		text width=2.2cm, minimum height=1cm, inner sep=3pt
	},
	flowarrow/.style={-{Stealth[length=2.2mm]}, primaryblue, thick},
	spokearrow/.style={{Stealth[length=1.8mm]}-{Stealth[length=1.8mm]}, primaryblue, thick},
	ucnode/.style={
		ellipse, draw=primaryblue, thick, fill=ucfill,
		align=center, font=\scriptsize\sffamily, text width=2.55cm,
		minimum height=1cm, inner sep=1pt
	},
	mainucnode/.style={
		ellipse, draw=primaryblue, very thick, fill=hubfill,
		align=center, font=\small\sffamily\bfseries, text width=2.9cm,
		minimum height=1.3cm, inner sep=2pt
	},
	extucnode/.style={
		ellipse, draw=primaryblue, thick, dashed, fill=extendfill,
		align=center, font=\scriptsize\sffamily, text width=2.55cm,
		minimum height=1cm, inner sep=1pt
	},
	includeedge/.style={-{Stealth[length=2mm]}, dashed, primaryblue, thick},
	extendedge/.style={-{Stealth[length=2mm]}, dashed, orange!70!black, thick},
	actorline/.style={-, actorgray, thick}
}

% Actor (stick figure) macro
\newcommand{\actorfigure}[3]{
	% #1 = center coordinate name, #2 = x, #3 = y
	\coordinate (#1) at (#2,#3);
	\draw[actorgray,thick] (#2,#3+0.9) circle (0.22);
	\draw[actorgray,thick] (#2,#3+0.68) -- (#2,#3-0.15);
	\draw[actorgray,thick] (#2-0.32,#3+0.45) -- (#2+0.32,#3+0.45);
	\draw[actorgray,thick] (#2,#3-0.15) -- (#2-0.3,#3-0.75);
	\draw[actorgray,thick] (#2,#3-0.15) -- (#2+0.3,#3-0.75);
}

\begin{document}
	
	% ===================== TITLE PAGE =====================
	\begin{titlepage}
		\centering
		\vspace*{1.5cm}
		{\headingfont\Huge\bfseries\color{primaryblue} SOFTWARE REQUIREMENTS SPECIFICATION\par}
		\vspace{0.6cm}
		{\headingfont\Large\bfseries HỆ THỐNG CÁ NHÂN HÓA LỘ TRÌNH HỌC VÀ LUYỆN THI\\ĐÁNH GIÁ NĂNG LỰC TÍCH HỢP AI\par}
		\vspace{1.5cm}
		{\large Tài liệu: Thiết kế Screen Flow, Use Case Diagram và Bảng mô tả màn hình\\cho vai trò \textbf{Administrator} và \textbf{Teacher / Academic Manager}\par}
		\vspace{2cm}
		
		\begin{tikzpicture}
			\node[hubnode, minimum width=6cm, minimum height=2.2cm, text width=5.5cm] {\Large AI\\[2pt]\normalsize Personalized Learning \& Assessment Platform};
		\end{tikzpicture}
		
		\vfill
		{\large Phiên bản: 1.0 \quad | \quad Ngày: \today\par}
	\end{titlepage}
	
	\tableofcontents
	\clearpage
	
	% =====================================================================
	============================
	\section{Sơ đồ Use Case chi tiết theo từng chức năng}
	% =====================================================================
	
	Mỗi sơ đồ dưới đây mô tả chi tiết một Use Case chính (tương ứng một nhóm chức năng),
	được phân rã thành các Use Case con thông qua quan hệ \textbf{«include»}
	(luôn thực hiện — mũi tên nét đứt đi từ Use Case chính đến Use Case con) và quan hệ
	\textbf{«extend»} (thực hiện có điều kiện/tùy chọn — mũi tên nét đứt đi
	từ Use Case mở rộng về Use Case chính). Do toàn bộ các quan hệ include/extend cũng như quan hệ
	association của Actor đều quy tụ về một điểm chung là Use Case chính, các đường trong mỗi sơ đồ
	\textbf{không giao/cắt nhau}.
	
	\subsection{1. Đăng nhập}
	\begin{figure}[h!]
		\centering
		\resizebox{\textwidth}{!}{%
			\begin{tikzpicture}
				
				\actorfigure{actlogin0}{-1.6}{0}
				\node[below=0.3cm of actlogin0, font=\scriptsize\sffamily\bfseries, align=center] {Người dùng\\(Admin/Teacher)};
				
				\node[mainucnode] (mainlogin) at (5.2,0) {Đăng nhập};
				\draw[actorline] (actlogin0) -- (mainlogin);
				
				\pgfmathsetmacro{\nsub}{4}
				
				\draw[draw=primaryblue, thick, rounded corners=6pt] (2.9,-2.73) rectangle (13.6,3.28);
				\node[font=\scriptsize\sffamily\bfseries, primaryblue] at (8.25,3.58) {Hệ thống — Đăng nhập};
				
				\node[ucnode] (inclogin0) at (10.4,1.88) {Nhập tài khoản\\và mật khẩu};
				\draw[includeedge] (mainlogin) -- node[midway, font=\tiny, fill=white, inner sep=1pt, align=center]{«include»} (inclogin0);
				\node[ucnode] (inclogin1) at (10.4,0.62) {Xác thực thông\\tin đăng nhập};
				\draw[includeedge] (mainlogin) -- node[midway, font=\tiny, fill=white, inner sep=1pt, align=center]{«include»} (inclogin1);
				\node[extucnode] (extlogin0) at (10.4,-0.62) {Quên mật khẩu};
				\draw[extendedge] (extlogin0) -- node[midway, font=\tiny, fill=white, inner sep=1pt, align=center]{«extend»} (mainlogin);
				\node[extucnode] (extlogin1) at (10.4,-1.88) {Xác thực hai lớp\\(2FA)};
				\draw[extendedge] (extlogin1) -- node[midway, font=\tiny, fill=white, inner sep=1pt, align=center]{«extend»} (mainlogin);
				
		\end{tikzpicture}}
		\caption{UC chi tiết – Đăng nhập}
	\end{figure}
	
	
	\subsection{2. Quản lý người dùng}
	\begin{figure}[h!]
		\centering
		\resizebox{\textwidth}{!}{%
			\begin{tikzpicture}
				
				\actorfigure{actusermgmt0}{-1.6}{0}
				\node[below=0.3cm of actusermgmt0, font=\scriptsize\sffamily\bfseries, align=center] {Administrator};
				
				\node[mainucnode] (mainusermgmt) at (5.2,0) {Quản lý\\người dùng};
				\draw[actorline] (actusermgmt0) -- (mainusermgmt);
				
				\pgfmathsetmacro{\nsub}{6}
				
				\draw[draw=primaryblue, thick, rounded corners=6pt] (2.9,-3.98) rectangle (13.6,4.53);
				\node[font=\scriptsize\sffamily\bfseries, primaryblue] at (8.25,4.83) {Hệ thống — Quản lý người dùng};
				
				\node[ucnode] (incusermgmt0) at (10.4,3.12) {Thêm người dùng};
				\draw[includeedge] (mainusermgmt) -- node[midway, font=\tiny, fill=white, inner sep=1pt, align=center]{«include»} (incusermgmt0);
				\node[ucnode] (incusermgmt1) at (10.4,1.88) {Sửa thông tin\\người dùng};
				\draw[includeedge] (mainusermgmt) -- node[midway, font=\tiny, fill=white, inner sep=1pt, align=center]{«include»} (incusermgmt1);
				\node[ucnode] (incusermgmt2) at (10.4,0.62) {Khóa / Mở\\tài khoản};
				\draw[includeedge] (mainusermgmt) -- node[midway, font=\tiny, fill=white, inner sep=1pt, align=center]{«include»} (incusermgmt2);
				\node[ucnode] (incusermgmt3) at (10.4,-0.62) {Tìm kiếm, lọc\\người dùng};
				\draw[includeedge] (mainusermgmt) -- node[midway, font=\tiny, fill=white, inner sep=1pt, align=center]{«include»} (incusermgmt3);
				\node[ucnode] (incusermgmt4) at (10.4,-1.88) {Gán vai trò};
				\draw[includeedge] (mainusermgmt) -- node[midway, font=\tiny, fill=white, inner sep=1pt, align=center]{«include»} (incusermgmt4);
				\node[extucnode] (extusermgmt0) at (10.4,-3.12) {Xuất danh sách\\người dùng};
				\draw[extendedge] (extusermgmt0) -- node[midway, font=\tiny, fill=white, inner sep=1pt, align=center]{«extend»} (mainusermgmt);
				
		\end{tikzpicture}}
		\caption{UC chi tiết – Quản lý người dùng}
	\end{figure}
	
	\clearpage
	
	\subsection{3. Quản lý vai trò}
	\begin{figure}[h!]
		\centering
		\resizebox{\textwidth}{!}{%
			\begin{tikzpicture}
				
				\actorfigure{actrolemgmt0}{-1.6}{0}
				\node[below=0.3cm of actrolemgmt0, font=\scriptsize\sffamily\bfseries, align=center] {Administrator};
				
				\node[mainucnode] (mainrolemgmt) at (5.2,0) {Quản lý vai trò\\và quyền};
				\draw[actorline] (actrolemgmt0) -- (mainrolemgmt);
				
				\pgfmathsetmacro{\nsub}{5}
				
				\draw[draw=primaryblue, thick, rounded corners=6pt] (2.9,-3.35) rectangle (13.6,3.90);
				\node[font=\scriptsize\sffamily\bfseries, primaryblue] at (8.25,4.20) {Hệ thống — Quản lý vai trò và quyền};
				
				\node[ucnode] (incrolemgmt0) at (10.4,2.50) {Tạo vai trò};
				\draw[includeedge] (mainrolemgmt) -- node[midway, font=\tiny, fill=white, inner sep=1pt, align=center]{«include»} (incrolemgmt0);
				\node[ucnode] (incrolemgmt1) at (10.4,1.25) {Sửa vai trò};
				\draw[includeedge] (mainrolemgmt) -- node[midway, font=\tiny, fill=white, inner sep=1pt, align=center]{«include»} (incrolemgmt1);
				\node[ucnode] (incrolemgmt2) at (10.4,0.00) {Xóa vai trò};
				\draw[includeedge] (mainrolemgmt) -- node[midway, font=\tiny, fill=white, inner sep=1pt, align=center]{«include»} (incrolemgmt2);
				\node[ucnode] (incrolemgmt3) at (10.4,-1.25) {Cấu hình ma\\trận quyền};
				\draw[includeedge] (mainrolemgmt) -- node[midway, font=\tiny, fill=white, inner sep=1pt, align=center]{«include»} (incrolemgmt3);
				\node[extucnode] (extrolemgmt0) at (10.4,-2.50) {Gán vai trò\\cho người dùng};
				\draw[extendedge] (extrolemgmt0) -- node[midway, font=\tiny, fill=white, inner sep=1pt, align=center]{«extend»} (mainrolemgmt);
				
		\end{tikzpicture}}
		\caption{UC chi tiết – Quản lý vai trò và quyền}
	\end{figure}
	
	
	\subsection{4. Quản lý môn học}
	\begin{figure}[h!]
		\centering
		\resizebox{\textwidth}{!}{%
			\begin{tikzpicture}
				
				\actorfigure{actsubjectmgmt0}{-1.6}{2.4}
				\node[below=0.3cm of actsubjectmgmt0, font=\scriptsize\sffamily\bfseries, align=center] {Administrator};
				\actorfigure{actsubjectmgmt1}{-1.6}{-2.4}
				\node[below=0.3cm of actsubjectmgmt1, font=\scriptsize\sffamily\bfseries, align=center] {Teacher};
				
				\node[mainucnode] (mainsubjectmgmt) at (5.2,0) {Quản lý\\môn học};
				\draw[actorline] (actsubjectmgmt0) -- (mainsubjectmgmt);
				\draw[actorline] (actsubjectmgmt1) -- (mainsubjectmgmt);
				
				\pgfmathsetmacro{\nsub}{5}
				
				\draw[draw=primaryblue, thick, rounded corners=6pt] (2.9,-3.35) rectangle (13.6,3.90);
				\node[font=\scriptsize\sffamily\bfseries, primaryblue] at (8.25,4.20) {Hệ thống — Quản lý môn học};
				
				\node[ucnode] (incsubjectmgmt0) at (10.4,2.50) {Thêm môn học};
				\draw[includeedge] (mainsubjectmgmt) -- node[midway, font=\tiny, fill=white, inner sep=1pt, align=center]{«include»} (incsubjectmgmt0);
				\node[ucnode] (incsubjectmgmt1) at (10.4,1.25) {Sửa môn học};
				\draw[includeedge] (mainsubjectmgmt) -- node[midway, font=\tiny, fill=white, inner sep=1pt, align=center]{«include»} (incsubjectmgmt1);
				\node[ucnode] (incsubjectmgmt2) at (10.4,0.00) {Ngừng / Kích hoạt\\môn học};
				\draw[includeedge] (mainsubjectmgmt) -- node[midway, font=\tiny, fill=white, inner sep=1pt, align=center]{«include»} (incsubjectmgmt2);
				\node[ucnode] (incsubjectmgmt3) at (10.4,-1.25) {Tìm kiếm\\môn học};
				\draw[includeedge] (mainsubjectmgmt) -- node[midway, font=\tiny, fill=white, inner sep=1pt, align=center]{«include»} (incsubjectmgmt3);
				\node[extucnode] (extsubjectmgmt0) at (10.4,-2.50) {Liên kết khung\\năng lực};
				\draw[extendedge] (extsubjectmgmt0) -- node[midway, font=\tiny, fill=white, inner sep=1pt, align=center]{«extend»} (mainsubjectmgmt);
				
		\end{tikzpicture}}
		\caption{UC chi tiết – Quản lý môn học}
	\end{figure}
	
	\clearpage
	
	\subsection{5. Quản lý khung năng lực}
	\begin{figure}[h!]
		\centering
		\resizebox{\textwidth}{!}{%
			\begin{tikzpicture}
				
				\actorfigure{actcompetencymgmt0}{-1.6}{2.4}
				\node[below=0.3cm of actcompetencymgmt0, font=\scriptsize\sffamily\bfseries, align=center] {Administrator};
				\actorfigure{actcompetencymgmt1}{-1.6}{-2.4}
				\node[below=0.3cm of actcompetencymgmt1, font=\scriptsize\sffamily\bfseries, align=center] {Teacher};
				
				\node[mainucnode] (maincompetencymgmt) at (5.2,0) {Quản lý khung\\năng lực};
				\draw[actorline] (actcompetencymgmt0) -- (maincompetencymgmt);
				\draw[actorline] (actcompetencymgmt1) -- (maincompetencymgmt);
				
				\pgfmathsetmacro{\nsub}{5}
				
				\draw[draw=primaryblue, thick, rounded corners=6pt] (2.9,-3.35) rectangle (13.6,3.90);
				\node[font=\scriptsize\sffamily\bfseries, primaryblue] at (8.25,4.20) {Hệ thống — Quản lý khung năng lực};
				
				\node[ucnode] (inccompetencymgmt0) at (10.4,2.50) {Tạo khung\\năng lực};
				\draw[includeedge] (maincompetencymgmt) -- node[midway, font=\tiny, fill=white, inner sep=1pt, align=center]{«include»} (inccompetencymgmt0);
				\node[ucnode] (inccompetencymgmt1) at (10.4,1.25) {Thêm tiêu chí\\năng lực};
				\draw[includeedge] (maincompetencymgmt) -- node[midway, font=\tiny, fill=white, inner sep=1pt, align=center]{«include»} (inccompetencymgmt1);
				\node[ucnode] (inccompetencymgmt2) at (10.4,0.00) {Thiết lập cấp độ\\năng lực};
				\draw[includeedge] (maincompetencymgmt) -- node[midway, font=\tiny, fill=white, inner sep=1pt, align=center]{«include»} (inccompetencymgmt2);
				\node[ucnode] (inccompetencymgmt3) at (10.4,-1.25) {Sửa / Xóa\\tiêu chí};
				\draw[includeedge] (maincompetencymgmt) -- node[midway, font=\tiny, fill=white, inner sep=1pt, align=center]{«include»} (inccompetencymgmt3);
				\node[extucnode] (extcompetencymgmt0) at (10.4,-2.50) {Liên kết\\môn học};
				\draw[extendedge] (extcompetencymgmt0) -- node[midway, font=\tiny, fill=white, inner sep=1pt, align=center]{«extend»} (maincompetencymgmt);
				
		\end{tikzpicture}}
		\caption{UC chi tiết – Quản lý khung năng lực}
	\end{figure}
	
	
	\subsection{6. Quản lý tài liệu}
	\begin{figure}[h!]
		\centering
		\resizebox{\textwidth}{!}{%
			\begin{tikzpicture}
				
				\actorfigure{actmaterialmgmt0}{-1.6}{2.4}
				\node[below=0.3cm of actmaterialmgmt0, font=\scriptsize\sffamily\bfseries, align=center] {Administrator};
				\actorfigure{actmaterialmgmt1}{-1.6}{-2.4}
				\node[below=0.3cm of actmaterialmgmt1, font=\scriptsize\sffamily\bfseries, align=center] {Teacher};
				
				\node[mainucnode] (mainmaterialmgmt) at (5.2,0) {Quản lý tài liệu\\học tập};
				\draw[actorline] (actmaterialmgmt0) -- (mainmaterialmgmt);
				\draw[actorline] (actmaterialmgmt1) -- (mainmaterialmgmt);
				
				\pgfmathsetmacro{\nsub}{5}
				
				\draw[draw=primaryblue, thick, rounded corners=6pt] (2.9,-3.35) rectangle (13.6,3.90);
				\node[font=\scriptsize\sffamily\bfseries, primaryblue] at (8.25,4.20) {Hệ thống — Quản lý tài liệu học tập};
				
				\node[ucnode] (incmaterialmgmt0) at (10.4,2.50) {Tải lên\\tài liệu};
				\draw[includeedge] (mainmaterialmgmt) -- node[midway, font=\tiny, fill=white, inner sep=1pt, align=center]{«include»} (incmaterialmgmt0);
				\node[ucnode] (incmaterialmgmt1) at (10.4,1.25) {Phân loại\\tài liệu};
				\draw[includeedge] (mainmaterialmgmt) -- node[midway, font=\tiny, fill=white, inner sep=1pt, align=center]{«include»} (incmaterialmgmt1);
				\node[ucnode] (incmaterialmgmt2) at (10.4,0.00) {Xóa tài liệu};
				\draw[includeedge] (mainmaterialmgmt) -- node[midway, font=\tiny, fill=white, inner sep=1pt, align=center]{«include»} (incmaterialmgmt2);
				\node[ucnode] (incmaterialmgmt3) at (10.4,-1.25) {Tìm kiếm\\tài liệu};
				\draw[includeedge] (mainmaterialmgmt) -- node[midway, font=\tiny, fill=white, inner sep=1pt, align=center]{«include»} (incmaterialmgmt3);
				\node[extucnode] (extmaterialmgmt0) at (10.4,-2.50) {Duyệt tài liệu};
				\draw[extendedge] (extmaterialmgmt0) -- node[midway, font=\tiny, fill=white, inner sep=1pt, align=center]{«extend»} (mainmaterialmgmt);
				
		\end{tikzpicture}}
		\caption{UC chi tiết – Quản lý tài liệu học tập}
	\end{figure}
	
	\clearpage
	
	\subsection{7. Quản lý ngân hàng câu hỏi}
	\begin{figure}[h!]
		\centering
		\resizebox{\textwidth}{!}{%
			\begin{tikzpicture}
				
				\actorfigure{actquestionbankmgmt0}{-1.6}{2.4}
				\node[below=0.3cm of actquestionbankmgmt0, font=\scriptsize\sffamily\bfseries, align=center] {Administrator};
				\actorfigure{actquestionbankmgmt1}{-1.6}{-2.4}
				\node[below=0.3cm of actquestionbankmgmt1, font=\scriptsize\sffamily\bfseries, align=center] {Teacher};
				
				\node[mainucnode] (mainquestionbankmgmt) at (5.2,0) {Quản lý ngân\\hàng câu hỏi};
				\draw[actorline] (actquestionbankmgmt0) -- (mainquestionbankmgmt);
				\draw[actorline] (actquestionbankmgmt1) -- (mainquestionbankmgmt);
				
				\pgfmathsetmacro{\nsub}{7}
				
				\draw[draw=primaryblue, thick, rounded corners=6pt] (2.9,-4.60) rectangle (13.6,5.15);
				\node[font=\scriptsize\sffamily\bfseries, primaryblue] at (8.25,5.45) {Hệ thống — Quản lý ngân hàng câu hỏi};
				
				\node[ucnode] (incquestionbankmgmt0) at (10.4,3.75) {Thêm câu hỏi};
				\draw[includeedge] (mainquestionbankmgmt) -- node[midway, font=\tiny, fill=white, inner sep=1pt, align=center]{«include»} (incquestionbankmgmt0);
				\node[ucnode] (incquestionbankmgmt1) at (10.4,2.50) {Sửa câu hỏi};
				\draw[includeedge] (mainquestionbankmgmt) -- node[midway, font=\tiny, fill=white, inner sep=1pt, align=center]{«include»} (incquestionbankmgmt1);
				\node[ucnode] (incquestionbankmgmt2) at (10.4,1.25) {Xóa câu hỏi};
				\draw[includeedge] (mainquestionbankmgmt) -- node[midway, font=\tiny, fill=white, inner sep=1pt, align=center]{«include»} (incquestionbankmgmt2);
				\node[ucnode] (incquestionbankmgmt3) at (10.4,0.00) {Phân loại\\độ khó};
				\draw[includeedge] (mainquestionbankmgmt) -- node[midway, font=\tiny, fill=white, inner sep=1pt, align=center]{«include»} (incquestionbankmgmt3);
				\node[ucnode] (incquestionbankmgmt4) at (10.4,-1.25) {Gắn thẻ\\năng lực};
				\draw[includeedge] (mainquestionbankmgmt) -- node[midway, font=\tiny, fill=white, inner sep=1pt, align=center]{«include»} (incquestionbankmgmt4);
				\node[extucnode] (extquestionbankmgmt0) at (10.4,-2.50) {Nhập câu hỏi\\từ file};
				\draw[extendedge] (extquestionbankmgmt0) -- node[midway, font=\tiny, fill=white, inner sep=1pt, align=center]{«extend»} (mainquestionbankmgmt);
				\node[extucnode] (extquestionbankmgmt1) at (10.4,-3.75) {Duyệt câu hỏi};
				\draw[extendedge] (extquestionbankmgmt1) -- node[midway, font=\tiny, fill=white, inner sep=1pt, align=center]{«extend»} (mainquestionbankmgmt);
				
		\end{tikzpicture}}
		\caption{UC chi tiết – Quản lý ngân hàng câu hỏi}
	\end{figure}
	
	
	\subsection{8. Quản lý đề thi}
	\begin{figure}[h!]
		\centering
		\resizebox{\textwidth}{!}{%
			\begin{tikzpicture}
				
				\actorfigure{actexammgmt0}{-1.6}{2.4}
				\node[below=0.3cm of actexammgmt0, font=\scriptsize\sffamily\bfseries, align=center] {Administrator};
				\actorfigure{actexammgmt1}{-1.6}{-2.4}
				\node[below=0.3cm of actexammgmt1, font=\scriptsize\sffamily\bfseries, align=center] {Teacher};
				
				\node[mainucnode] (mainexammgmt) at (5.2,0) {Quản lý đề thi\\và bài đánh giá};
				\draw[actorline] (actexammgmt0) -- (mainexammgmt);
				\draw[actorline] (actexammgmt1) -- (mainexammgmt);
				
				\pgfmathsetmacro{\nsub}{6}
				
				\draw[draw=primaryblue, thick, rounded corners=6pt] (2.9,-3.98) rectangle (13.6,4.53);
				\node[font=\scriptsize\sffamily\bfseries, primaryblue] at (8.25,4.83) {Hệ thống — Quản lý đề thi và bài đánh giá};
				
				\node[ucnode] (incexammgmt0) at (10.4,3.12) {Tạo đề thi\\thủ công};
				\draw[includeedge] (mainexammgmt) -- node[midway, font=\tiny, fill=white, inner sep=1pt, align=center]{«include»} (incexammgmt0);
				\node[ucnode] (incexammgmt1) at (10.4,1.88) {Cấu hình thời\\gian làm bài};
				\draw[includeedge] (mainexammgmt) -- node[midway, font=\tiny, fill=white, inner sep=1pt, align=center]{«include»} (incexammgmt1);
				\node[ucnode] (incexammgmt2) at (10.4,0.62) {Gán đối tượng\\dự thi};
				\draw[includeedge] (mainexammgmt) -- node[midway, font=\tiny, fill=white, inner sep=1pt, align=center]{«include»} (incexammgmt2);
				\node[ucnode] (incexammgmt3) at (10.4,-0.62) {Xuất bản\\đề thi};
				\draw[includeedge] (mainexammgmt) -- node[midway, font=\tiny, fill=white, inner sep=1pt, align=center]{«include»} (incexammgmt3);
				\node[extucnode] (extexammgmt0) at (10.4,-1.88) {Tạo đề tự động\\theo ma trận (AI)};
				\draw[extendedge] (extexammgmt0) -- node[midway, font=\tiny, fill=white, inner sep=1pt, align=center]{«extend»} (mainexammgmt);
				\node[extucnode] (extexammgmt1) at (10.4,-3.12) {Thu hồi\\đề thi};
				\draw[extendedge] (extexammgmt1) -- node[midway, font=\tiny, fill=white, inner sep=1pt, align=center]{«extend»} (mainexammgmt);
				
		\end{tikzpicture}}
		\caption{UC chi tiết – Quản lý đề thi và bài đánh giá}
	\end{figure}
	
	\clearpage
	
	\subsection{9. Quản lý thông báo}
	\begin{figure}[h!]
		\centering
		\resizebox{\textwidth}{!}{%
			\begin{tikzpicture}
				
				\actorfigure{actnotifmgmt0}{-1.6}{2.4}
				\node[below=0.3cm of actnotifmgmt0, font=\scriptsize\sffamily\bfseries, align=center] {Administrator};
				\actorfigure{actnotifmgmt1}{-1.6}{-2.4}
				\node[below=0.3cm of actnotifmgmt1, font=\scriptsize\sffamily\bfseries, align=center] {Teacher};
				
				\node[mainucnode] (mainnotifmgmt) at (5.2,0) {Quản lý\\thông báo};
				\draw[actorline] (actnotifmgmt0) -- (mainnotifmgmt);
				\draw[actorline] (actnotifmgmt1) -- (mainnotifmgmt);
				
				\pgfmathsetmacro{\nsub}{5}
				
				\draw[draw=primaryblue, thick, rounded corners=6pt] (2.9,-3.35) rectangle (13.6,3.90);
				\node[font=\scriptsize\sffamily\bfseries, primaryblue] at (8.25,4.20) {Hệ thống — Quản lý thông báo};
				
				\node[ucnode] (incnotifmgmt0) at (10.4,2.50) {Soạn\\thông báo};
				\draw[includeedge] (mainnotifmgmt) -- node[midway, font=\tiny, fill=white, inner sep=1pt, align=center]{«include»} (incnotifmgmt0);
				\node[ucnode] (incnotifmgmt1) at (10.4,1.25) {Chọn đối tượng\\nhận};
				\draw[includeedge] (mainnotifmgmt) -- node[midway, font=\tiny, fill=white, inner sep=1pt, align=center]{«include»} (incnotifmgmt1);
				\node[ucnode] (incnotifmgmt2) at (10.4,0.00) {Gửi thông báo};
				\draw[includeedge] (mainnotifmgmt) -- node[midway, font=\tiny, fill=white, inner sep=1pt, align=center]{«include»} (incnotifmgmt2);
				\node[ucnode] (incnotifmgmt3) at (10.4,-1.25) {Sửa / Xóa\\thông báo};
				\draw[includeedge] (mainnotifmgmt) -- node[midway, font=\tiny, fill=white, inner sep=1pt, align=center]{«include»} (incnotifmgmt3);
				\node[extucnode] (extnotifmgmt0) at (10.4,-2.50) {Lên lịch gửi};
				\draw[extendedge] (extnotifmgmt0) -- node[midway, font=\tiny, fill=white, inner sep=1pt, align=center]{«extend»} (mainnotifmgmt);
				
		\end{tikzpicture}}
		\caption{UC chi tiết – Quản lý thông báo}
	\end{figure}
	
	
	\subsection{10. Báo cáo}
	\begin{figure}[h!]
		\centering
		\resizebox{\textwidth}{!}{%
			\begin{tikzpicture}
				
				\actorfigure{actreport0}{-1.6}{2.4}
				\node[below=0.3cm of actreport0, font=\scriptsize\sffamily\bfseries, align=center] {Administrator};
				\actorfigure{actreport1}{-1.6}{-2.4}
				\node[below=0.3cm of actreport1, font=\scriptsize\sffamily\bfseries, align=center] {Teacher};
				
				\node[mainucnode] (mainreport) at (5.2,0) {Xem báo cáo\\và thống kê};
				\draw[actorline] (actreport0) -- (mainreport);
				\draw[actorline] (actreport1) -- (mainreport);
				
				\pgfmathsetmacro{\nsub}{4}
				
				\draw[draw=primaryblue, thick, rounded corners=6pt] (2.9,-2.73) rectangle (13.6,3.28);
				\node[font=\scriptsize\sffamily\bfseries, primaryblue] at (8.25,3.58) {Hệ thống — Báo cáo và thống kê};
				
				\node[ucnode] (increport0) at (10.4,1.88) {Lọc theo\\tiêu chí};
				\draw[includeedge] (mainreport) -- node[midway, font=\tiny, fill=white, inner sep=1pt, align=center]{«include»} (increport0);
				\node[ucnode] (increport1) at (10.4,0.62) {Xem biểu đồ\\thống kê};
				\draw[includeedge] (mainreport) -- node[midway, font=\tiny, fill=white, inner sep=1pt, align=center]{«include»} (increport1);
				\node[ucnode] (increport2) at (10.4,-0.62) {Xem chi tiết\\báo cáo};
				\draw[includeedge] (mainreport) -- node[midway, font=\tiny, fill=white, inner sep=1pt, align=center]{«include»} (increport2);
				\node[extucnode] (extreport0) at (10.4,-1.88) {Xuất báo cáo\\(PDF/Excel)};
				\draw[extendedge] (extreport0) -- node[midway, font=\tiny, fill=white, inner sep=1pt, align=center]{«extend»} (mainreport);
				
		\end{tikzpicture}}
		\caption{UC chi tiết – Báo cáo và thống kê}
	\end{figure}
	
	\clearpage
	
	\subsection{11. Theo dõi tiến độ}
	\begin{figure}[h!]
		\centering
		\resizebox{\textwidth}{!}{%
			\begin{tikzpicture}
				
				\actorfigure{actprogress0}{-1.6}{0}
				\node[below=0.3cm of actprogress0, font=\scriptsize\sffamily\bfseries, align=center] {Teacher};
				
				\node[mainucnode] (mainprogress) at (5.2,0) {Theo dõi tiến độ\\và báo cáo năng lực};
				\draw[actorline] (actprogress0) -- (mainprogress);
				
				\pgfmathsetmacro{\nsub}{4}
				
				\draw[draw=primaryblue, thick, rounded corners=6pt] (2.9,-2.73) rectangle (13.6,3.28);
				\node[font=\scriptsize\sffamily\bfseries, primaryblue] at (8.25,3.58) {Hệ thống — Theo dõi tiến độ và báo cáo năng lực};
				
				\node[ucnode] (incprogress0) at (10.4,1.88) {Xem tiến độ\\theo lớp};
				\draw[includeedge] (mainprogress) -- node[midway, font=\tiny, fill=white, inner sep=1pt, align=center]{«include»} (incprogress0);
				\node[ucnode] (incprogress1) at (10.4,0.62) {Xem xu hướng\\năng lực};
				\draw[includeedge] (mainprogress) -- node[midway, font=\tiny, fill=white, inner sep=1pt, align=center]{«include»} (incprogress1);
				\node[ucnode] (incprogress2) at (10.4,-0.62) {Lọc theo lớp /\\môn / kỳ};
				\draw[includeedge] (mainprogress) -- node[midway, font=\tiny, fill=white, inner sep=1pt, align=center]{«include»} (incprogress2);
				\node[extucnode] (extprogress0) at (10.4,-1.88) {Xuất báo cáo\\tiến độ};
				\draw[extendedge] (extprogress0) -- node[midway, font=\tiny, fill=white, inner sep=1pt, align=center]{«extend»} (mainprogress);
				
		\end{tikzpicture}}
		\caption{UC chi tiết – Theo dõi tiến độ và báo cáo năng lực}
	\end{figure}
	
	
	\subsection{12. Hồ sơ năng lực}
	\begin{figure}[h!]
		\centering
		\resizebox{\textwidth}{!}{%
			\begin{tikzpicture}
				
				\actorfigure{actcompetencyprofile0}{-1.6}{0}
				\node[below=0.3cm of actcompetencyprofile0, font=\scriptsize\sffamily\bfseries, align=center] {Teacher};
				
				\node[mainucnode] (maincompetencyprofile) at (5.2,0) {Xem hồ sơ\\năng lực học sinh};
				\draw[actorline] (actcompetencyprofile0) -- (maincompetencyprofile);
				
				\pgfmathsetmacro{\nsub}{4}
				
				\draw[draw=primaryblue, thick, rounded corners=6pt] (2.9,-2.73) rectangle (13.6,3.28);
				\node[font=\scriptsize\sffamily\bfseries, primaryblue] at (8.25,3.58) {Hệ thống — Xem hồ sơ năng lực học sinh};
				
				\node[ucnode] (inccompetencyprofile0) at (10.4,1.88) {Xem bản đồ\\năng lực};
				\draw[includeedge] (maincompetencyprofile) -- node[midway, font=\tiny, fill=white, inner sep=1pt, align=center]{«include»} (inccompetencyprofile0);
				\node[ucnode] (inccompetencyprofile1) at (10.4,0.62) {Xem chi tiết\\theo tiêu chí};
				\draw[includeedge] (maincompetencyprofile) -- node[midway, font=\tiny, fill=white, inner sep=1pt, align=center]{«include»} (inccompetencyprofile1);
				\node[ucnode] (inccompetencyprofile2) at (10.4,-0.62) {So sánh\\theo thời gian};
				\draw[includeedge] (maincompetencyprofile) -- node[midway, font=\tiny, fill=white, inner sep=1pt, align=center]{«include»} (inccompetencyprofile2);
				\node[extucnode] (extcompetencyprofile0) at (10.4,-1.88) {Xem theo\\môn học};
				\draw[extendedge] (extcompetencyprofile0) -- node[midway, font=\tiny, fill=white, inner sep=1pt, align=center]{«extend»} (maincompetencyprofile);
				
		\end{tikzpicture}}
		\caption{UC chi tiết – Xem hồ sơ năng lực học sinh}
	\end{figure}
	
	\clearpage
	
	% =====================================================================
	\section{Bảng mô tả màn hình}
	% =====================================================================
	
	Ký hiệu cột: \textbf{Nhóm CN} = Tên nhóm chức năng; \textbf{Màn hình liên kết} = các màn hình có
	thể điều hướng đến/từ màn hình hiện tại; \textbf{UC liên quan} = Use Case liên quan (xem Mục 2).
	
	% ---------------------------------------------------------------
	\subsection{Bảng mô tả màn hình – Administrator}
	% ---------------------------------------------------------------
	
	\begin{landscape}
		\footnotesize
		\begin{longtable}{|>{\raggedright\arraybackslash}p{2.1cm}|>{\raggedright\arraybackslash}p{2.3cm}|>{\raggedright\arraybackslash}p{3.4cm}|>{\raggedright\arraybackslash}p{3.6cm}|>{\raggedright\arraybackslash}p{3.6cm}|>{\raggedright\arraybackslash}p{2.6cm}|>{\raggedright\arraybackslash}p{2.6cm}|}
			\caption{Bảng mô tả màn hình – Administrator}\\
			\hline
			\rowcolor{lightblue}
			\textbf{Nhóm CN} & \textbf{Tên màn hình} & \textbf{Mục đích màn hình} & \textbf{Dữ liệu hiển thị} & \textbf{Thao tác chính} & \textbf{Màn hình liên kết} & \textbf{UC liên quan} \\
			\hline
			\endfirsthead
			\hline
			\rowcolor{lightblue}
			\textbf{Nhóm CN} & \textbf{Tên màn hình} & \textbf{Mục đích màn hình} & \textbf{Dữ liệu hiển thị} & \textbf{Thao tác chính} & \textbf{Màn hình liên kết} & \textbf{UC liên quan} \\
			\hline
			\endhead
			\hline
			\endfoot
			
			Xác thực &
			Đăng nhập &
			Xác thực tài khoản Admin để truy cập hệ thống &
			Form usern/mật khẩu, thông báo lỗi, liên kết quên mật khẩu &
			Nhập thông tin, đăng nhập, quên mật khẩu, xác thực 2 lớp (nếu có) &
			Dashboard &
			Đăng nhập \\
			\hline
			
			Tổng quan &
			Dashboard &
			Cung cấp cái nhìn tổng quan số liệu vận hành toàn hệ thống &
			Số liệu người dùng, số đề thi, số bài đánh giá, cảnh báo hệ thống, biểu đồ nhanh &
			Xem số liệu tổng hợp, điều hướng nhanh tới các module, lọc theo thời gian &
			Tất cả các màn hình chức năng &
			Xem Dashboard tổng quan \\
			\hline
			
			Người dùng &
			Quản lý người dùng &
			Quản lý toàn bộ tài khoản người dùng trong hệ thống &
			Danh sách người dùng, trạng thái, vai trò, thông tin liên hệ &
			Thêm/sửa/xóa, khóa/mở tài khoản, gán vai trò, tìm kiếm, lọc, xuất danh sách &
			Dashboard, Quản lý vai trò và quyền &
			Quản lý người dùng \\
			\hline
			
			Phân quyền &
			Quản lý vai trò và quyền &
			Định nghĩa vai trò và gán quyền truy cập chức năng &
			Danh sách vai trò, ma trận quyền theo chức năng/module &
			Tạo/sửa/xóa vai trò, cấu hình quyền, gán vai trò cho người dùng &
			Dashboard, Quản lý người dùng &
			Quản lý vai trò và phân quyền \\
			\hline
			
			Học thuật &
			Quản lý môn học &
			Quản trị danh mục môn học của hệ thống &
			Danh sách môn học, mã môn, mô tả, trạng thái &
			Thêm/sửa/xóa môn học, kích hoạt/ngừng, tìm kiếm &
			Dashboard, Quản lý khung năng lực &
			Quản lý môn học \\
			\hline
			
			Học thuật &
			Quản lý khung năng lực &
			Xây dựng và quản trị khung năng lực (competency framework) theo môn học/cấp độ &
			Cây năng lực, mức độ, tiêu chí đánh giá, môn học liên kết &
			Thêm/sửa/xóa tiêu chí năng lực, thiết lập cấp độ, liên kết môn học &
			Dashboard, Quản lý môn học, Quản lý ngân hàng câu hỏi &
			Quản lý khung năng lực \\
			\hline
			
			Nội dung &
			Quản lý tài liệu học tập &
			Quản trị kho tài liệu học tập dùng chung toàn hệ thống &
			Danh sách tài liệu, loại tệp, môn học, người tải lên, trạng thái duyệt &
			Tải lên/xóa/duyệt tài liệu, phân loại theo môn học/năng lực, tìm kiếm &
			Dashboard, Quản lý môn học &
			Quản lý tài liệu học tập \\
			\hline
			
			Nội dung &
			Quản lý ngân hàng câu hỏi &
			Quản trị kho câu hỏi dùng cho đề thi và bài đánh giá năng lực &
			Danh sách câu hỏi, độ khó, dạng câu hỏi, môn học, năng lực liên kết &
			Thêm/sửa/xóa câu hỏi, phân loại độ khó, gắn thẻ năng lực, nhập từ file, duyệt &
			Dashboard, Quản lý khung năng lực, Quản lý đề thi và bài đánh giá &
			Quản lý ngân hàng câu hỏi \\
			\hline
			
			Đánh giá &
			Quản lý đề thi và bài đánh giá &
			Tạo lập, cấu hình và phát hành đề thi/bài đánh giá năng lực &
			Danh sách đề thi, cấu trúc đề, thời lượng, trạng thái phát hành &
			Tạo đề (thủ công/tự động theo ma trận), cấu hình thời gian, xuất bản, thu hồi &
			Dashboard, Quản lý ngân hàng câu hỏi, Báo cáo và thống kê &
			Quản lý đề thi và bài đánh giá \\
			\hline
			
			Vận hành &
			Quản lý thông báo &
			Soạn thảo và gửi thông báo hệ thống tới người dùng &
			Danh sách thông báo, đối tượng nhận, thời gian gửi, trạng thái &
			Soạn/sửa/xóa thông báo, chọn đối tượng nhận, lên lịch gửi &
			Dashboard &
			Quản lý thông báo \\
			\hline
			
			Hệ thống &
			Cấu hình và giám sát hệ thống &
			Cấu hình tham số hệ thống và giám sát tình trạng vận hành, nhật ký &
			Tham số cấu hình, log hệ thống, trạng thái dịch vụ/AI engine, cảnh báo &
			Cập nhật cấu hình, xem log, khởi động lại dịch vụ, xử lý cảnh báo &
			Dashboard &
			Cấu hình và giám sát hệ thống \\
			\hline
			
			Báo cáo &
			Báo cáo và thống kê &
			Tổng hợp báo cáo và thống kê toàn hệ thống phục vụ ra quyết định &
			Biểu đồ, bảng số liệu theo người dùng/môn học/kỳ đánh giá &
			Lọc theo tiêu chí, xem chi tiết, xuất báo cáo (PDF/Excel) &
			Dashboard, Quản lý đề thi và bài đánh giá &
			Xem báo cáo và thống kê \\
			\hline
			
		\end{longtable}
	\end{landscape}
	
	\clearpage
	
	% ---------------------------------------------------------------
	\subsection{Bảng mô tả màn hình – Teacher / Academic Manager}
	% ---------------------------------------------------------------
	
	\begin{landscape}
		\footnotesize
		\begin{longtable}{|>{\raggedright\arraybackslash}p{2.1cm}|>{\raggedright\arraybackslash}p{2.3cm}|>{\raggedright\arraybackslash}p{3.4cm}|>{\raggedright\arraybackslash}p{3.6cm}|>{\raggedright\arraybackslash}p{3.6cm}|>{\raggedright\arraybackslash}p{2.6cm}|>{\raggedright\arraybackslash}p{2.6cm}|}
			\caption{Bảng mô tả màn hình – Teacher / Academic Manager}\\
			\hline
			\rowcolor{lightblue}
			\textbf{Nhóm CN} & \textbf{Tên màn hình} & \textbf{Mục đích màn hình} & \textbf{Dữ liệu hiển thị} & \textbf{Thao tác chính} & \textbf{Màn hình liên kết} & \textbf{UC liên quan} \\
			\hline
			\endfirsthead
			\hline
			\rowcolor{lightblue}
			\textbf{Nhóm CN} & \textbf{Tên màn hình} & \textbf{Mục đích màn hình} & \textbf{Dữ liệu hiển thị} & \textbf{Thao tác chính} & \textbf{Màn hình liên kết} & \textbf{UC liên quan} \\
			\hline
			\endhead
			\hline
			\endfoot
			
			Xác thực &
			Đăng nhập &
			Xác thực tài khoản Teacher/Academic Manager &
			Form username/mật khẩu, liên kết quên mật khẩu &
			Nhập thông tin, đăng nhập, quên mật khẩu &
			Dashboard &
			Đăng nhập \\
			\hline
			
			Tổng quan &
			Dashboard &
			Tổng quan lớp học, môn học phụ trách và tình trạng học sinh &
			Số môn phụ trách, số học sinh, cảnh báo tiến độ, thông báo mới &
			Xem số liệu tổng hợp, điều hướng nhanh tới các module &
			Tất cả các màn hình chức năng &
			Xem Dashboard tổng quan \\
			\hline
			
			Học thuật &
			Quản lý môn học và khung năng lực &
			Quản trị môn học phụ trách và khung năng lực liên quan &
			Danh sách môn học, cây năng lực, tiêu chí đánh giá &
			Sửa/đề xuất môn học, cập nhật khung năng lực, gán tiêu chí &
			Dashboard, Quản lý tài liệu học tập &
			Quản lý môn học và khung năng lực \\
			\hline
			
			Nội dung &
			Quản lý tài liệu học tập &
			Biên soạn và quản lý tài liệu học tập cho môn học phụ trách &
			Danh sách tài liệu, loại tệp, môn học, trạng thái duyệt &
			Tải lên/sửa/xóa tài liệu, phân loại theo năng lực, gửi duyệt &
			Dashboard, Quản lý môn học và khung năng lực &
			Quản lý tài liệu học tập \\
			\hline
			
			Nội dung &
			Quản lý ngân hàng câu hỏi &
			Biên soạn câu hỏi phục vụ đề thi và bài đánh giá &
			Danh sách câu hỏi, độ khó, dạng câu hỏi, năng lực liên kết &
			Thêm/sửa/xóa câu hỏi, gắn thẻ năng lực, gửi duyệt &
			Dashboard, Quản lý đề thi và bài đánh giá &
			Quản lý ngân hàng câu hỏi \\
			\hline
			
			Đánh giá &
			Quản lý đề thi và bài đánh giá &
			Tạo và tổ chức đề thi/bài đánh giá cho lớp phụ trách &
			Danh sách đề thi, cấu trúc đề, thời lượng, đối tượng dự thi &
			Tạo đề, gán lớp/học sinh, cấu hình thời gian, xuất bản &
			Dashboard, Quản lý ngân hàng câu hỏi, Xem kết quả học tập &
			Quản lý đề thi và bài đánh giá \\
			\hline
			
			Kết quả &
			Xem kết quả học tập &
			Theo dõi kết quả làm bài/đánh giá của học sinh &
			Điểm số, thời gian làm bài, câu đúng/sai theo từng học sinh &
			Xem chi tiết, lọc theo lớp/kỳ đánh giá, xuất báo cáo &
			Dashboard, Quản lý đề thi và bài đánh giá, Xem hồ sơ năng lực học sinh &
			Xem kết quả học tập \\
			\hline
			
			Hồ sơ HS &
			Xem hồ sơ năng lực học sinh &
			Xem chi tiết năng lực hiện tại của từng học sinh &
			Bản đồ năng lực, mức độ đạt được theo từng tiêu chí &
			Xem chi tiết theo học sinh/môn học, so sánh theo thời gian &
			Dashboard, Xem kết quả học tập, Xem lộ trình học &
			Xem hồ sơ năng lực học sinh \\
			\hline
			
			Lộ trình &
			Xem lộ trình học &
			Theo dõi lộ trình học cá nhân hóa do AI đề xuất cho học sinh &
			Danh sách nội dung/bài tập theo lộ trình, tiến độ hoàn thành &
			Xem chi tiết lộ trình, gợi ý điều chỉnh, phản hồi cho hệ thống AI &
			Dashboard, Xem hồ sơ năng lực học sinh, Theo dõi tiến độ và báo cáo năng lực &
			Xem lộ trình học \\
			\hline
			
			Báo cáo &
			Theo dõi tiến độ và báo cáo năng lực &
			Tổng hợp tiến độ học tập và năng lực theo lớp/nhóm học sinh &
			Biểu đồ tiến độ, thống kê năng lực theo lớp, xu hướng theo thời gian &
			Lọc theo lớp/môn/kỳ, xem chi tiết, xuất báo cáo &
			Dashboard, Xem lộ trình học, Xem hồ sơ năng lực học sinh &
			Theo dõi tiến độ và báo cáo năng lực \\
			\hline
			
			Vận hành &
			Quản lý thông báo và hồ sơ cá nhân &
			Gửi thông báo tới học sinh/lớp và quản trị hồ sơ cá nhân &
			Danh sách thông báo, thông tin cá nhân, cài đặt tài khoản &
			Soạn/gửi thông báo, cập nhật hồ sơ cá nhân, đổi mật khẩu &
			Dashboard &
			Quản lý thông báo và hồ sơ cá nhân \\
			\hline
			
		\end{longtable}
	\end{landscape}
	
\end{document}